\scanpage{29}\mspage{28}
% @ top, false start
\deleted{Alors}

Ainsi

% @ top, proposition lines 1–3
\uline{Proposition} On a construit un Inf de $F$ et $G$ pour $\ll$, on \deleted{\unsure{y|l}}l'appelle la ``figure intersection'' $F.G$.

% @ upper middle, warning lines 1–2
Attention, ce \uline{n'est pas} \inserted{\unsure{en général}} la figure $F\cap G$, \deleted{\unsure{intersection|inter.}} en se rappelant $F,\deleted{F}G\subset\unsure{\mathfrak P|\mathscr P|\mathbb P}(L)$, \deleted{$G\subset\mathfrak P(L)$}

% @ middle, paragraph lines 1–3
Nous n'allons considérer cette $F\cap G$ que dans le cas où $F$ et $G$ compatibles, mais alors les deux coïncident~!

% @ middle, paragraph lines 1–4
Donc la bonne notion d'intersection est le $F.G$. En fait, ça marche pour une famille quelconque de figures, et
% @ middle lower, support formula
\[
 \left|\bigwedge_{i\in I} F_i\right|=\bigcap_{i\in I}|F_i|\quad .
\]

% @ lower part, paragraph lines 1–4
Cette propriété d'existence d'un Inf pour la relation $\ll$ \unsure{s'apparente|l'apparente} \deleted{\unsure{la priorité|la propriété|la présentation}} à une propriété de ``\uline{\unsure{modération|médiation}}''.

% @ foot, cancelled statement lines 1–3
\deleted{\uline{Propriété} Si $F'\preceq F$ \deleted{\unsure{, et des}}, alors $F'.G\preceq F.G$ (\unsure{Immédiat|Immédit|immédite} sur supports)~: Subdivision \uline{induite} sur un raffinement.}
